% UML class diagram: sim::Environment holds the pluggable physics models --
% gravity, atmosphere, geoid -- behind string-keyed registries; WindModel is
% compiled but dormant (zero call sites). Bare tikzpicture; the whitepaper
% wraps it in a figure and \resizebox{\linewidth}. Names verified against
% notes/environment.md (sim/Environment.h and friends).
\begin{tikzpicture}[node distance=8mm and 9mm]

  % ---- container (top center) ----
  \node[umlclass,text width=46mm,anchor=north] (env) at (74mm,0mm) {%
    \textbf{Environment}\\
    {\scriptsize non-copyable; one per app, owned by QtRocket}
    \nodepart{two}{\scriptsize\ttfamily
      - gravityModels : map<string,\\
      \hspace*{2.5mm}shared\_ptr<GravityModel>>\\
      - atmosphereModels : map<string,\\
      \hspace*{2.5mm}shared\_ptr<AtmosphericModel>>\\
      - gravityModel : string {\sffamily\color{qrGray}(selection)}\\
      - atmosphereModel : string\\
      - geoidModel :\\
      \hspace*{2.5mm}shared\_ptr<GeoidModel>}
    \nodepart{three}{\scriptsize\ttfamily
      + setGravityModel(name : string)\\
      + setAtmosphereModel(name : string)\\
      + getGravityModel()\\
      \hspace*{2.5mm}: shared\_ptr<GravityModel>\\
      + getAtmosphericModel()\\
      \hspace*{2.5mm}: shared\_ptr<AtmosphericModel>\\
      + getGeoidModel()\\
      \hspace*{2.5mm}: shared\_ptr<GeoidModel>}};

  % ---- registry note (right of the container) ----
  \node[umlnote,text width=40mm,right=5mm of env] (regnote) {%
    model registries are string-keyed:\\[1pt]
    {\ttfamily "Constant Gravity",\\ "Spherical Gravity"}\\[1pt]
    {\ttfamily "Constant Atmosphere",\\ "US Standard 1976", "Vacuum"}\\[2pt]
    {\color{qrGray}unknown name $\Rightarrow$ silent no-op; each
     \texttt{set*Model} call re-creates a fresh instance}};
  \draw[qrGray,line width=0.4pt,dash pattern=on 2pt off 2pt]
    (regnote.west) -- (env.east);

  % ---- abstract bases (one per column) ----
  \node[umlabstract,text width=42mm,anchor=north]
      (grav) at ($(env.south)+(-51.4mm,-12mm)$) {%
    {\scriptsize\guillemotleft abstract\guillemotright}\\
    \textit{\textbf{GravityModel}}
    \nodepart{two}
    \nodepart{three}{\scriptsize\ttfamily
      + getAccel(x, y, z) : Vector3\\
      + getAccel(Vector3) : Vector3\\
      {\sffamily\color{qrGray}accel in m/s$^2$, local launch frame}}};

  \node[umlabstract,text width=36mm,anchor=north]
      (geoid) at ($(env.south)+(0mm,-12mm)$) {%
    {\scriptsize\guillemotleft abstract\guillemotright}\\
    \textit{\textbf{GeoidModel}}
    \nodepart{two}
    \nodepart{three}{\scriptsize\ttfamily
      + getGroundLevel(lat, lon)\\
      \hspace*{2.5mm}: double\\
      {\sffamily\color{qrGray}ground distance from Earth's center, m}}};

  \node[umlabstract,text width=42mm,anchor=north]
      (atmo) at ($(env.south)+(51.4mm,-12mm)$) {%
    {\scriptsize\guillemotleft abstract\guillemotright}\\
    \textit{\textbf{AtmosphericModel}}
    \nodepart{two}
    \nodepart{three}{\scriptsize\ttfamily
      + getDensity(altitude)\\
      + getPressure(altitude)\\
      + getTemperature(altitude)\\
      + getSpeedOfSound(altitude)\\
      + getDynamicViscosity(altitude)\\
      {\sffamily\color{qrGray}all pure virtual; altitude in m above
       MSL, all return double}}};

  % ---- gravity implementations (left column) ----
  \node[umlclass,text width=42mm,below=6mm of grav] (cgrav) {%
    \textbf{ConstantGravityModel}
    \nodepart{two}
    \nodepart{three}{\scriptsize\ttfamily
      + getAccel(x, y, z) : Vector3\\
      {\sffamily\color{qrGray}uniform $(0,\,0,\,-g_0)$;
       $g_0$ = 9.80665 m/s$^2$}}};

  \node[umlclass,text width=42mm,below=6mm of cgrav] (sgrav) {%
    \textbf{SphericalGravityModel}
    \nodepart{two}{\scriptsize\ttfamily
      - geoid : shared\_ptr<GeoidModel>\\
      - groundLevel : double\\
      \hspace*{2.5mm}{\sffamily\color{qrGray}geoid radius, cached at ctor}}
    \nodepart{three}{\scriptsize\ttfamily
      + SphericalGravityModel(geoid)\\
      + getAccel(x, y, z) : Vector3\\
      {\sffamily\color{qrGray}$\vec a = -GM\,\vec r_{geo}/r^3$, geocentric
       via $z + \mathrm{groundLevel}$; finite at the pad}}};

  % ---- geoid implementation (middle column) ----
  \node[umlclass,text width=36mm,below=6mm of geoid] (sgeoid) {%
    \textbf{SphericalGeoidModel}
    \nodepart{two}
    \nodepart{three}{\scriptsize\ttfamily
      + getGroundLevel(lat, lon)\\
      {\sffamily\color{qrGray}always WGS84 mean radius\\
       6\,371\,008.8 m; ignores lat/lon}}};

  % ---- atmosphere implementations (right column) ----
  \node[umlclass,text width=42mm,below=6mm of atmo] (catmo) {%
    \textbf{ConstantAtmosphere}
    \nodepart{two}
    \nodepart{three}{\scriptsize\sffamily\color{qrGray}
      all five getters return sea-level\\
      constants at any altitude\\
      (1.225 kg/m$^3$, 101\,325 Pa, 288.15 K)}};

  \node[umlclass,text width=42mm,below=6mm of catmo] (usatmo) {%
    \textbf{USStandardAtmosphere}
    \nodepart{two}{\scriptsize\ttfamily
      - temperatureLapseRate : Bin\\
      - standardTemperature : Bin\\
      - standardPressure : Bin\\
      - standardDensity : Bin\\
      {\sffamily\color{qrGray}static 7-layer 1976 tables}}
    \nodepart{three}{\scriptsize\sffamily\color{qrGray}
      overrides all five getters;\\
      analytic lapse-rate / isothermal\\
      formulas; throws below 0 m}};

  \node[umlclass,text width=42mm,below=6mm of usatmo] (vacatmo) {%
    \textbf{VacuumAtmosphere}
    \nodepart{two}
    \nodepart{three}{\scriptsize\sffamily\color{qrGray}
      all five getters return 0.0;\\
      Mach callers must guard $a = 0$}};

  % ---- wind: a seam with no consumer (middle column, standalone) ----
  \node[umlclass,text width=36mm,below=11mm of sgeoid] (wind) {%
    \textbf{WindModel}
    \nodepart{two}
    \nodepart{three}{\scriptsize\ttfamily
      + getWindSpeed(x, y, z)\\
      \hspace*{2.5mm}: Vector3\\
      {\sffamily\color{qrGray}virtual, not pure ---\\
       base returns $(0,\,0,\,0)$}}};

  \node[umlnote,text width=36mm,below=5mm of wind] (windnote) {%
    dormant --- zero call sites\\[1pt]
    {\color{qrGray}compiled into \texttt{sim/} but never instantiated;
     not held by Environment. drag treats the airmass as at rest.}};
  \draw[qrGray,line width=0.4pt,dash pattern=on 2pt off 2pt]
    (windnote.north) -- (wind.south);

  % ---- inheritance (arrows point at the base) ----
  \draw[inherit] (cgrav.north) -- (grav.south);
  \draw[inherit] (sgrav.west) -- ++(-4mm,0) |- (grav.west);
  \draw[inherit] (sgeoid.north) -- (geoid.south);
  \draw[inherit] (catmo.north) -- (atmo.south);
  \draw[inherit] (usatmo.east) -- ++(4mm,0) |- (atmo.east);
  \draw[inherit] (vacatmo.east) -- ++(4mm,0) |- (atmo.east);

  % ---- aggregation: Environment holds the models by shared_ptr; the
  % multiplicities are the registry sizes (2 gravity / 3 atmosphere keys) ----
  \draw[aggregate] ([xshift=-16mm]env.south) -- ++(0,-5mm)
    -| node[umllbl,pos=0.25,above]{\ttfamily gravityModels}
       node[mult,pos=0.92,right=1pt]{2} (grav.north);
  \draw[aggregate] (env.south) -- (geoid.north)
    node[umllbl,midway,right=1pt]{\ttfamily geoidModel}
    node[mult,pos=0.8,left=1pt]{1};
  \draw[aggregate] ([xshift=16mm]env.south) -- ++(0,-5mm)
    -| node[umllbl,pos=0.25,above]{\ttfamily atmosphereModels}
       node[mult,pos=0.92,left=1pt]{3} (atmo.north);

  % ---- SphericalGravityModel takes the geoid at construction ----
  \draw[assoc] (sgrav.east) -- ++(4mm,0) |-
    node[umllbl,pos=0.72,above,xshift=-2.5mm]{\ttfamily geoid} (geoid.west);

\end{tikzpicture}
